%% kapa: exam=1 round=20 year=2009 period=1 form=A subject=economics subject_ko=경제원론 seq=12 number=52
%% source: Q-Net 감정평가사 기출문제 (cst00309) artlSeq=5024178 file=extracted/1008/5024178/2009년 제20회 감정평가사 1차시험 기출문제/1교시.hwp
%% conversion: kordoc 4.24.51 -> tools/gen1.py (auto); review=proofread-sonnet; flags=none
\begin{question}{52}
A국가는 상대적으로 노동력이 풍부하고 B국가는 상대적으로 자본이 풍부하다. 헥셔-올린 정리에 입각할 때 경제적 교류가 전혀 없던 두 국가 간에 자유무역이 이루어진다면, 무역 이전과 비교하여 A국가의 임금과 이자율의 변화는?
\begin{choices}
  \item 임금은 상승하고 이자율은 하락할 것이다.
  \item 임금은 하락하고 이자율은 상승할 것이다.
  \item 임금과 이자율 모두 하락할 것이다.
  \item 임금과 이자율 모두 상승할 것이다.
  \item 임금과 이자율 모두 불변일 것이다.
\end{choices}
\end{question}
\answer{1}
